% GENERATED FILE. Edit canonical lessons, then run npm run generate:books.
% canonical-source-hash: fnv1a64:62db1ec33061fc98
% canonical-lessons: TE-C329-dongatanam, TE-C329-corabadu, TE-C329-baddalu-kottu, TE-C329-ettukupovu, TE-C329-pasigattu

\chapter{A Break-In}
\label{ch:te-a-break-in}
\hlchaptermodality{329}

\begin{chapteropening}
\textbf{By the end of this chapter:} \emph{I can say a theft, to break in, to break open, to make off with and to sense.}

Complete the last lesson of chapter 329: I can say a theft, to break in, to break open, to make off with and to sense.
\end{chapteropening}

\section[\texorpdfstring{\te{దొంగతనం} (doṅgatanaṁ)}{దొంగతనం (doṅgatanaṁ)}]{\te{దొంగతనం} (doṅgatanaṁ) — a theft; a burglary}

\label{lesson:TE-C329-dongatanam}



\begin{quote}
Before the new one: say the Telugu for the next day, then the Telugu for to recall.
\end{quote}

\subsection*{You'll want to know: \te{దొంగతనం}}
\textbf{\te{దొంగతనం}} — \emph{doṅgatanaṁ} — \textquotedblleft{}a theft; a burglary\textquotedblright{}.

\begin{quote}
\te{నిన్న} \te{రాత్రి} \te{మా} \te{వీధిలో} \te{దొంగతనం} \te{జరిగింది}.
\end{quote}

\begin{quote}
\emph{ninna rātri mā vīdhilō doṅgatanaṁ jarigindi.}
\end{quote}

\begin{quote}
Last night there was a theft in our street.
\end{quote}

\te{దొంగ}, a thief, with -\te{తనం}, the ending that makes a noun of a state or a deed.

\begin{cousinweb}
\textbf{\te{దొంగతనం}} (\emph{doṅgatanaṁ}) is native Dravidian: \emph{doṅga}, \textquotedblleft{}a thief\textquotedblright{} (DEDR 2982), with the Telugu noun ending \emph{-tanaṁ}.
\end{cousinweb}

\begin{grammarlens}[title={What you've built}]
One more word for telling what happened.
\end{grammarlens}

\subsection*{Your turn}
\begin{itemize}
\raggedright
  \item \emph{Say it:} \emph{doṅgatanaṁ}
  \item \emph{Say it:} \emph{doṅgatanaṁ}, once more
  \item \emph{Say it:} \emph{ninna rātri mā vīdhilō doṅgatanaṁ jarigindi.} — Last night there was a theft in our street.
  \item \emph{Recall:} say the Telugu for the next day, then the Telugu for to recall, then say \emph{doṅgatanaṁ} again
\end{itemize}

\begin{tcolorbox}[breakable,colback=teal!4,colframe=teal!35!black,title={Before you move on}]
What does \te{దొంగతనం} mean? (\textquotedblleft{}A theft; a burglary\textquotedblright{}.) Say it once more.
\end{tcolorbox}
\section[\texorpdfstring{\te{చొరబడు} (corabaḍu)}{చొరబడు (corabaḍu)}]{\te{చొరబడు} (corabaḍu) — to break in; to intrude}

\label{lesson:TE-C329-corabadu}



\begin{quote}
Before the new one: say the Telugu for to recall, then the Telugu for a theft.
\end{quote}

\subsection*{You'll want to know: \te{చొరబడు}}
\textbf{\te{చొరబడు}} — \emph{corabaḍu} — \textquotedblleft{}to break in; to intrude\textquotedblright{}.

\begin{quote}
\te{దొంగలు} \te{కిటికీ} \te{నుంచి} \te{ఇంట్లోకి} \te{చొరబడ్డారు}.
\end{quote}

\begin{quote}
\emph{doṅgalu kiṭikī nuñci iṇṭlōki corabaḍḍāru.}
\end{quote}

\begin{quote}
The thieves broke into the house through the window.
\end{quote}

Past (I …): \te{చొరబడ్డాను}, \emph{corabaḍḍānu}.

The place you enter takes -\te{లోకి}, into: \te{ఇంట్లోకి} \te{చొరబడు}.

\begin{cousinweb}
\textbf{\te{చొరబడు}} (\emph{corabaḍu}) is a native compound: the stem \emph{cor-} of \emph{coccu}, \textquotedblleft{}to enter, penetrate\textquotedblright{} (DEDR 2778), with \emph{paḍu}, the helping verb of things that befall.
\end{cousinweb}

\begin{grammarlens}[title={What you've built}]
One more word for telling what happened.
\end{grammarlens}

\subsection*{Your turn}
\begin{itemize}
\raggedright
  \item \emph{Say it:} \emph{corabaḍu}
  \item \emph{Say it:} \emph{corabaḍu}, once more
  \item \emph{Say it:} \emph{doṅgalu kiṭikī nuñci iṇṭlōki corabaḍḍāru.} — The thieves broke into the house through the window.
  \item \emph{Recall:} say the Telugu for to recall, then the Telugu for a theft, then say \emph{corabaḍu} again
\end{itemize}

\begin{tcolorbox}[breakable,colback=teal!4,colframe=teal!35!black,title={Before you move on}]
What does \te{చొరబడు} mean? (\textquotedblleft{}To break in; to intrude\textquotedblright{}.) Say it once more.
\end{tcolorbox}
\section[\texorpdfstring{\te{బద్దలు} \te{కొట్టు} (baddalu koṭṭu)}{బద్దలు కొట్టు (baddalu koṭṭu)}]{\te{బద్దలు} \te{కొట్టు} (baddalu koṭṭu) — to break open; to smash}

\label{lesson:TE-C329-baddalu-kottu}



\begin{quote}
Before the new one: say the Telugu for a theft, then the Telugu for to break in.
\end{quote}

\subsection*{You'll want to know: \te{బద్దలు} \te{కొట్టు}}
\textbf{\te{బద్దలు} \te{కొట్టు}} — \emph{baddalu koṭṭu} — \textquotedblleft{}to break open; to smash\textquotedblright{}.

\begin{quote}
\te{దొంగ} \te{తలుపు} \te{తాళం} \te{బద్దలు} \te{కొట్టాడు}.
\end{quote}

\begin{quote}
\emph{doṅga talupu tāḷaṁ baddalu koṭṭāḍu.}
\end{quote}

\begin{quote}
The thief smashed the lock on the door.
\end{quote}

Past (I …): \te{బద్దలు} \te{కొట్టాను}, \emph{baddalu koṭṭānu}.

\te{బద్దలు} means pieces: to beat something to pieces is to break it open.

\begin{cousinweb}
\textbf{\te{బద్దలు} \te{కొట్టు}} (\emph{baddalu koṭṭu}) is a Telugu phrase: \emph{koṭṭu}, \textquotedblleft{}to beat, strike\textquotedblright{}, is native Dravidian (DEDR 2063); the origin of \emph{baddalu}, \textquotedblleft{}pieces\textquotedblright{}, is not settled.
\end{cousinweb}

\begin{grammarlens}[title={What you've built}]
One more word for telling what happened.
\end{grammarlens}

\subsection*{Your turn}
\begin{itemize}
\raggedright
  \item \emph{Say it:} \emph{baddalu koṭṭu}
  \item \emph{Say it:} \emph{baddalu koṭṭu}, once more
  \item \emph{Say it:} \emph{doṅga talupu tāḷaṁ baddalu koṭṭāḍu.} — The thief smashed the lock on the door.
  \item \emph{Recall:} say the Telugu for a theft, then the Telugu for to break in, then say \emph{baddalu koṭṭu} again
\end{itemize}

\begin{tcolorbox}[breakable,colback=teal!4,colframe=teal!35!black,title={Before you move on}]
What does \te{బద్దలు} \te{కొట్టు} mean? (\textquotedblleft{}To break open; to smash\textquotedblright{}.) Say it once more.
\end{tcolorbox}
\section[\texorpdfstring{\te{ఎత్తుకుపోవు} (ettukupōvu)}{ఎత్తుకుపోవు (ettukupōvu)}]{\te{ఎత్తుకుపోవు} (ettukupōvu) — to make off with}

\label{lesson:TE-C329-ettukupovu}



\begin{quote}
Before the new one: say the Telugu for to break in, then the Telugu for to break open.
\end{quote}

\subsection*{You'll want to know: \te{ఎత్తుకుపోవు}}
\textbf{\te{ఎత్తుకుపోవు}} — \emph{ettukupōvu} — \textquotedblleft{}to make off with\textquotedblright{}.

\begin{quote}
\te{దొంగలు} \te{బంగారం} \te{అంతా} \te{ఎత్తుకుపోయారు}.
\end{quote}

\begin{quote}
\emph{doṅgalu baṅgāraṁ antā ettukupōyāru.}
\end{quote}

\begin{quote}
The thieves made off with all the gold.
\end{quote}

Past (I …): \te{ఎత్తుకుపోయాను}, \emph{ettukupōyānu}.

Three pieces in one: \te{ఎత్తు}, to lift, -\te{కు}, for oneself, and \te{పోవు}, to go.

\begin{cousinweb}
\textbf{\te{ఎత్తుకుపోవు}} (\emph{ettukupōvu}) is a native compound: \emph{ettu}, \textquotedblleft{}to lift up, take away\textquotedblright{} (DEDR 796), with the self-ending \emph{-ku} and \emph{pōvu}, \textquotedblleft{}to go\textquotedblright{} (DEDR 4572).
\end{cousinweb}

\begin{grammarlens}[title={What you've built}]
One more word for telling what happened.
\end{grammarlens}

\subsection*{Your turn}
\begin{itemize}
\raggedright
  \item \emph{Say it:} \emph{ettukupōvu}
  \item \emph{Say it:} \emph{ettukupōvu}, once more
  \item \emph{Say it:} \emph{doṅgalu baṅgāraṁ antā ettukupōyāru.} — The thieves made off with all the gold.
  \item \emph{Recall:} say the Telugu for to break in, then the Telugu for to break open, then say \emph{ettukupōvu} again
\end{itemize}

\begin{tcolorbox}[breakable,colback=teal!4,colframe=teal!35!black,title={Before you move on}]
What does \te{ఎత్తుకుపోవు} mean? (\textquotedblleft{}To make off with\textquotedblright{}.) Say it once more.
\end{tcolorbox}
\section[\texorpdfstring{\te{పసిగట్టు} (pasigaṭṭu)}{పసిగట్టు (pasigaṭṭu)}]{\te{పసిగట్టు} (pasigaṭṭu) — to sense; to detect}

\label{lesson:TE-C329-pasigattu}



\begin{quote}
Before the new one: say the Telugu for to break open, then the Telugu for to make off with.
\end{quote}

\subsection*{You'll want to know: \te{పసిగట్టు}}
\textbf{\te{పసిగట్టు}} — \emph{pasigaṭṭu} — \textquotedblleft{}to sense; to detect\textquotedblright{}.

\begin{quote}
\te{కుక్క} \te{దొంగను} \te{వెంటనే} \te{పసిగట్టింది}.
\end{quote}

\begin{quote}
\emph{kukka doṅganu veṇṭanē pasigaṭṭindi.}
\end{quote}

\begin{quote}
The dog sensed the thief at once.
\end{quote}

Past (I …): \te{పసిగట్టాను}, \emph{pasigaṭṭānu}.

A dog senses a stranger, and so does a sharp mind: \te{పసిగట్టు} works for both.

\begin{cousinweb}
\textbf{\te{పసిగట్టు}} (\emph{pasigaṭṭu}) is a Telugu compound whose first part is \emph{pasi}, \textquotedblleft{}smell, scent\textquotedblright{} (DEDR 3826); the origin of the whole word is not settled.
\end{cousinweb}

\begin{grammarlens}[title={What you've built}]
One more word for telling what happened.
\end{grammarlens}

\subsection*{Your turn}
\begin{itemize}
\raggedright
  \item \emph{Say it:} \emph{pasigaṭṭu}
  \item \emph{Say it:} \emph{pasigaṭṭu}, once more
  \item \emph{Say it:} \emph{kukka doṅganu veṇṭanē pasigaṭṭindi.} — The dog sensed the thief at once.
  \item \emph{Recall:} say the Telugu for to break open, then the Telugu for to make off with, then say \emph{pasigaṭṭu} again
\end{itemize}

\begin{tcolorbox}[breakable,colback=teal!4,colframe=teal!35!black,title={Before you move on}]
What does \te{పసిగట్టు} mean? (\textquotedblleft{}To sense; to detect\textquotedblright{}.) Say it once more.
\end{tcolorbox}
